\subsection{Закон Джоуля – Ленца}
%%%%%%%%%%%%%%%%%%%%%%%%%%%%%%%%%%%%%%%%%%%%%%%%%%%%
%%%%%%%%%%%%%%%%%%%%%%%%%%%%%%%%%%%%%%%%%%%%%%%%%%%%
%%%%%%%%%%%%%%%%%%%%%%%%%%%%%%%%%%%%%%%%%%%%%%%%%%%%
 \z{Электрическая плитка при силе тока $I = 5\,\text{А}$ за время $t = 3\,\text{мин}$ потребляет энергию $W = 1080\,\text{кДж}$. Рассчитать сопротивление плитки.}
%
{$R = \frac{Q}{I^2 t} = 240\,\text{Ом}$}
%%%%%%%%%%%%%%%%%%%%%%%%%%%%%%%%%%%%%%%%%%%%%%%%%%%%
 \z[{\begin{wrapfigure}[3]{r}{0.2\textwidth}%
		\vspace{-30pt}%
		\begin{flushright}%
 			\includegraphics[width=0.2\textwidth]{3.5.2}%
		\end{flushright}%
	\end{wrapfigure}%
	}]%
{Три лампы одинаковой мощности, рассчитанные на одно и то же напряжение, включены как показано на рисунке. Одинаков ли накал нитей ламп?}
%
{Максимальная мощность на левой лампочке}

%%%%%%%%%%%%%%%%%%%%%%%%%%%%%%%%%%%%%%%%%%%%%%%%%%%%
 \z[{\begin{wrapfigure}[3]{r}{0.2\textwidth}%
		\vspace{-20pt}%
		\begin{flushright}%
 			\includegraphics[width=0.2\textwidth]{3.5.3}%
		\end{flushright}%
	\end{wrapfigure}%
	}]%
{Три лампы, имеющие одинаковые сопротивления, включены в цепь так, как показано на рисунке. Во сколько раз мощность тока в лампе 1 больше мощности тока в лампе 2?}
%
{в 4 раза}

%%%%%%%%%%%%%%%%%%%%%%%%%%%%%%%%%%%%%%%%%%%%%%%%%%%%
 \z[{\begin{wrapfigure}[6]{r}{0.2\textwidth}%
		\vspace{-12pt}%
		\begin{flushright}%
 			\includegraphics[width=0.2\textwidth]{3.5.4}%
		\end{flushright}%
	\end{wrapfigure}%
	}]%
{В сеть напряжением $U = 120\,\text{В}$ параллельно включены две лампы: Л1 - мощностью $P = 300\,\text{Вт}$, рассчитанную на напряжение $120\,\text{В}$, и Л2 - 12-вольтная лампа, включенная последовательно с проводником c сопротивлением $R$. Лампа Л2 работает в номинальном режиме. Определить показания амперметров $A_1$ и $A$ и сопротивление проводника, если амперметр $A_2$ показывает силу тока $I = 2\,\text{А}$.}
%
{$I_1 = 2.5\,\text{А},\quad I_2 = 4.5\,\text{А},\quad R = 54\,\text{Ом}$}

%%%%%%%%%%%%%%%%%%%%%%%%%%%%%%%%%%%%%%%%%%%%%%%%%%%%
 \z{Лампа накаливания имеет сопротивление нити в нагретом состоянии $R = 300\,\text{Ом}$ и работает от сети напряжением $U = 127\,\text{В}$. На сколько градусов нагреется сосуд с водой за время $t = 5\,\text{мин}$, если в него поместить лампочку? Теплоемкость сосуда $C = 10\,\text{Дж/К}$, масса воды $m = 1.2\,\text{кг}$.}
%
{$\Delta T = \frac{U^2 t}{R(cm + C)} = 3.1\,\degree\text{C}$}

%%%%%%%%%%%%%%%%%%%%%%%%%%%%%%%%%%%%%%%%%%%%%%%%%%%%
 \z{К источнику тока с ЭДС $\emf = 3\,\text{В}$ и внутренним сопротивлением $r = 1\,\text{Ом}$ подключено сопротивление $R = 2\,\text{Ом}$. Найти полезную мощность источника тока, мощность источника и потери мощности в цепи.}
%
{$P_{\text{полез}} = \frac{\emf^2 R}{(R + r)^2} = 2\,\text{Вт};\quad P_{\text{ист}} = \frac{\emf^2}{R + r} = 3\,\text{Вт};\quad P_{\text{пот}} = \frac{\emf^2 r}{(R + r)^2} = 1\,\text{Вт}$}

%%%%%%%%%%%%%%%%%%%%%%%%%%%%%%%%%%%%%%%%%%%%%%%%%%%%
 \z[{\begin{wrapfigure}[5]{r}{0.25\textwidth}%
		\vspace{-15pt}%
		\begin{flushright}%
 			\includegraphics[width=0.25\textwidth]{3.5.7}%
		\end{flushright}%
	\end{wrapfigure}%
	}]%
{Электрическая цепь (рис.) подключена к сети постоянного напряжения. При изменении сопротивления переменного резистора $R$ на нём выделяется мощность $P_0 = 16\,\text{Вт}$ при токах $I_1 = 1\,\text{А}$ и $I_2 = 4\,\text{А}$. Определите наибольшую мощность $P_{\text{max}}$, которая может выделиться на резисторе $R$.}
%
{$P_{\text{max}} = P_0 \frac{(I_1 + I_2)^2}{4I_1 I_2} = 25\,\text{Вт}$}